\documentclass{article}

\usepackage{arxiv}
\usepackage[utf8]{inputenc}
\usepackage[T1]{fontenc}
\usepackage{hyperref}
\usepackage{url}
\usepackage{booktabs}
\usepackage{amsfonts}
\usepackage{nicefrac}
\usepackage{microtype}
\usepackage{graphicx}
\usepackage{amsmath}
\graphicspath{ {./images/} {./DAlambert/} {./} }


% override
\renewcommand{\refname}{Referencias}
\renewcommand{\contentsname}{Índice}
\usepackage{fancyhdr}
\fancyhead[R]{\today}
% -----------------------------------

\title{TP 2.1 - GENERADORES PSEUDOALEATORIOS}
\author{
  Manuel Ferraris \\
  Universidad Tecnológica Nacional - FRRO \\
  Zeballos 1341, S2000, Argentina \\
  \texttt{ferrarismanu@gmail.com} \\
  \And
  Juan Ignacio Fiele \\
  Universidad Tecnológica Nacional - FRRO \\
  Zeballos 1341, S2000, Argentina \\
  \texttt{juanfiele2002@gmail.com} \\
  \And
  María Belén Giannone \\
  Universidad Tecnológica Nacional - FRRO \\
  Zeballos 1341, S2000, Argentina \\
  \texttt{mariabelengiannone010@gmail.com} \\
  \And
  María Lucía Munné \\
  Universidad Tecnológica Nacional - FRRO \\
  Zeballos 1341, S2000, Argentina \\
  \texttt{marilumunne@gmail.com} \\
  \And
  Agustín Stringa \\
  Universidad Tecnológica Nacional - FRRO \\
  Zeballos 1341, S2000, Argentina \\
  \texttt{stringaagustin1@gmail.com} \\
}
\date{\today}
\begin{document}
\maketitle

\keywords{Python \and Numeros Aleatorios \and Generadores Pseudoaleatorios \and Simulación}

\begin{abstract}
\end{abstract}

\section{Introducción}

\section{Desarrollo Teórico}

\subsection{Generadores de numeros pseudoaleatorios implementados}


\subsection{Pruebas de aleatoriedad}
Las pruebas de aleatoriedad se aplican a los generadores para evaluar su calidad.
Si bien, una prueba simple puede realizarse mediante observación visual de ciertos gráficos, no son en absoluto suficientes para validar un generador.
Por ello, se aplican pruebas estadísticas formales que evalúan diferentes aspectos de la secuencia generada, como la uniformidad, la independencia y la ausencia de patrones predecibles.

Entre las pruebas estadísticas se pueden diferenciar:

\begin{itemize}
  \item \textbf{empíricas:} Evalúan estadísticas de sucesiones de números.
  \item \textbf{teóricas:} Se establecen las características de las sucesiones usando métodos de teoría de números con base en la regla de recurrencia que generó la sucesión.
\end{itemize}


Prueba de bondad de ajuste \textbf{chi cuadrado}
Esta prueba determina si una sola variable categórica se ajusta a una distribución teórica esperada.
Las hipótesis de la misma son:
\begin{itemize}
  \item $H_0$: La variable se ajusta a la distribución teórica \textit{F}.
  \item $H_1$: La variable no se ajusta a la distribución teórica \textit{F}.
\end{itemize}

Procedimiento:
\begin{list}{label}{spacing}
  \item [1.] Partir el soporte de F en c celdas que son exhaustivas y mutuamente
  \item [2.] Contar el número de observaciones $O_i$ que caen en cada celda i.
\end{list}

Para evaluar la calidad del generador de números pseudoaleatorios en el dominio $[0, 1)$, se realizó una partición del espacio muestral en $k=10$ sub-intervalos equiprobables de amplitud $0.1$. La hipótesis nula ($H_0$) establece que el generador sigue una distribución uniforme continua, lo que implica que la frecuencia esperada $E_i$ para cada intervalo $i$ es constante e igual a $n/k$, donde $n$ representa el tamaño de la muestra

n = 10000

% https://es.wikipedia.org/wiki/Prueba_%CF%87%C2%B2

\section{Resultados}

\begin{figure}[htbp]
  \centering
  \includegraphics[width=1\textwidth]{./images/estrategia_o_capital_f.png}
  \caption{Análisis integral de la simulación larga ($t=200, c=20$): Convergencia de métricas, distribución de frecuencias y validación del TCL.}
  \label{fig:sim_larga}
\end{figure}

\section{Discusión y Conclusiones}


\begin{thebibliography}{9}

  \bibitem{enunciado}
  Cátedra de Simulación. (2026). \textit{Enunciado Trabajo Práctico 1.1: Simulación de una Ruleta}. Universidad Tecnológica Nacional - Facultad Regional Rosario (UTN-FRRO).

  \bibitem{catedra}
  Cátedra de Probabilidad y Estadística. \textit{Tratamiento de Datos}. Universidad Tecnológica Nacional - Facultad Regional Rosario (UTN-FRRO). [2024].

  \bibitem{wikipedia_ruleta}
  Wikipedia. \textit{Ruleta}. Recuperado de: \url{https://es.wikipedia.org/wiki/Ruleta}.

  \bibitem{arxiv_template}
  Cornell University. \textit{arXiv Template for Preprints}.

  \bibitem{matplotlib}
  Matplotlib Development Team. (2024). \textit{Matplotlib API Reference Documentation}.
  Recuperado de: \url{https://matplotlib.org/stable/api/index.html}

  \bibitem{numpy}
  NumPy Developers. (2024). \textit{NumPy Reference Documentation}.
  Recuperado de: \url{https://numpy.org/doc/stable/reference/index.html}

  \bibitem{cpython_random}
  Python Software Foundation. \textit{Source code for random.py}. GitHub - CPython Repository.
  Recuperado de: \url{https://github.com/python/cpython/blob/main/Lib/random.py}

\end{thebibliography}

\end{document}